\chapter{Abstract} \begin{tabular}{@{} >{\bfseries}p{0.2\textwidth} p{0.77\textwidth} @{}} \MakeUppercase{Keywords} & clinical decision support; MIMIC--IV; automated ICD coding; rare disease prediction; synthetic data generation; NLP; sepsis-associated acute kidney injury (SA--AKI); Sepsis--3; KDIGO; early warning; risk stratification; leakage-aware modeling; calibration; SHAP explanations; decision-curve analysis; survival analysis; UMLS; RxNorm; SNOMED CT\\ \end{tabular} \par \begin{doublespacing}
This thesis investigates two complementary machine learning problems for clinical decision support, both grounded in the MIMIC--IV database.

\textbf{Problem~1: Automated ICD Code Prediction for Rare Diseases.} The long-tailed distribution of ICD-10-CM codes means that rare but clinically significant conditions are chronically under-represented in training data, leading to poor recall where coding errors carry the greatest consequence. This problem was investigated through an extensive literature review spanning rule-based systems, deep learning architectures (CNNs, RNNs, Transformers), curriculum learning, and synthetic data generation. A formal problem definition, thorough dataset exploration of MIMIC--IV discharge summaries (364,627 patients; 19,440 unique ICD-10 codes), and preliminary modeling were completed. A KNN baseline achieved a Macro F1-score of 0.12, confirming the difficulty of the task, and a rare-disease co-occurrence pipeline was developed using the ORDO ontology and SNOMED~CT. Full model training and synthetic data generation were not completed due to computational infrastructure costs and remain as future work.

\textbf{Problem~2: Early Mortality Risk Stratification for SA--AKI.} At the end of the first ICU day (T0--T24), patients with sepsis-associated acute kidney injury who are at the highest risk of in-hospital death must be reliably identified. Previous SA--AKI models, while often reporting strong discrimination, have been found to suffer from temporal leakage, global fitting of preprocessing steps, and a lack of calibration. To overcome these limitations, a leakage-averse, TRIPOD-aligned framework was developed. An ontology-integrated data warehouse was built using PostgreSQL to harmonize MIMIC--IV with medical terminologies (UMLS, RxNorm, SNOMED~CT, ATC). The SA--AKI cohort (n=10,036) was rigorously defined using Sepsis--3 and KDIGO criteria, with all predictors strictly confined to the T0--T24 window. A one-time 80/20 stratified split ensured that all data-dependent steps---including a novel approach to creating and selecting significant missingness indicators---were fitted only on training data. The methodology encompassed three stages: (i)~exploratory data analysis to identify clinical phenotypes; (ii)~supervised mortality prediction, from which a LightGBM classifier was selected, tuned with Optuna, and calibrated using isotonic regression; and (iii)~survival analysis using Cox Proportional Hazards and a custom DeepSurv model. On the held-out test set, the final calibrated LightGBM model achieved an AUROC of \textbf{0.749} and an F1-score of 0.522. The survival models yielded C-indexes of 0.693 (Cox) and 0.688 (DeepSurv). SHAP analysis confirmed that predictions were driven by established clinical risk factors.

The primary contribution of this work is a unified clinical decision support framework addressing both structured prediction (ICD coding) and temporal risk stratification (SA--AKI mortality), with a reproducible, timing-faithful pipeline that delivers calibrated, interpretable outputs suitable for bedside decision support.
\end{doublespacing}
